% =====================================================================
% Appendix: Inductor design and eddy-current measurement data
% Sources: Excel Models/Ansys to Loss.xlsx, Water and LM models.xlsx
% (Comparison sheet), inductor Rac.xlsx, Experimental Measurement.xlsx
% (Galinstan - Inductor, LCR Measurements)
% =====================================================================
\chapter{Inductor and Eddy-Current Data}
\label{app:inductor-data}

This appendix gives the data behind Experiment~2: the inductor's design values
and measured parameters, the material properties that govern eddy-current loss
at \SI{85}{\kilo\hertz}, all AC resistance measurements, and the flow and coolant
temperatures recorded while the inductor was energised.

\section{Inductor Parameters}

\Cref{tab:app-inductor} compares the final design with the built inductor. The
winding loss estimated in the design spreadsheet, which assumed that the AC
resistance of the 800-strand Litz wire is twice its DC resistance, agrees with the
loss calculated from the measured AC resistance. The ANSYS sweeps over air-gap
length and number of turns, from which the design was chosen, are held in the
project file \texttt{Ansys to Loss.xlsx}.

\begin{table}[htbp]
  \centering
  \caption{Design and measured parameters of the inductor.}
  \label{tab:app-inductor}
  \small
  \begin{tabular}{@{}lcc@{}}
    \toprule
    Parameter & Design (ANSYS) & Measured \\
    \midrule
    Turns, air gap & 8, \SI{0.6}{\milli\metre} & -- \\
    Inductance (\si{\micro\henry}) & 92.5 & 92.1 \\
    Core loss at \SI{12}{\ampere} rms (\si{\watt}) & 30.81 & -- \\
    Winding AC resistance (\si{\milli\ohm}) & -- & 83 \\
    Winding loss at \SI{12}{\ampere} rms (\si{\watt}) & 11.93 & 11.95 \\
    Total loss (\si{\watt}) & 42.74 & 42.76\tnote{a} \\
    \bottomrule
  \end{tabular}
  \floatnote{Core at \SI{65}{\celsius}; AC resistance at \SI{85}{\kilo\hertz}
    and low excitation. \tnote{a}Design core loss plus measured winding loss.}
\end{table}

\section{Skin Depth and Eddy-Current Loss of the Materials}

\Cref{tab:app-skin} gives the skin depth $\delta = 1/\sqrt{\pi f \mu_0 \sigma}$
of the materials near the inductor at \SI{85}{\kilo\hertz}, and the eddy-current
loss each develops as a plate in the same alternating field, from the initial
ANSYS simulation described in \cref{sec:inductor-design}. Galinstan's lower
conductivity gives a skin depth about three to four times that of aluminium or
copper. In a plate thicker than its skin depth, the loss falls as the
conductivity rises, so Galinstan dissipates 3.2 times as much as aluminium and
3.9 times as much as copper. Water, by contrast, is effectively transparent to
the field.

\begin{table}[htbp]
  \centering
  \caption{Skin depth and simulated eddy-current loss at \SI{85}{\kilo\hertz}.}
  \label{tab:app-skin}
  \small
  \begin{tabular}{@{}lccc@{}}
    \toprule
    Material & $\sigma$ & $\delta$ & Plate loss relative \\
             & (\si{\siemens\per\metre}) & (\si{\milli\metre}) & to copper \\
    \midrule
    Copper    & \num{5.96e7} & 0.22 & 1.0 \\
    Aluminium & \num{3.77e7} & 0.28 & 1.2 \\
    Galinstan & \num{3.46e6} & 0.93 & 3.9 \\
    Water     & \num{8e-4}   & \num{6.1e4} & -- \\
    \bottomrule
  \end{tabular}
  \floatnote{Plate losses from an ANSYS Maxwell tutorial geometry, used only to
    compare the materials.}
\end{table}

\section{AC Resistance Measurements}

All AC resistances were measured with the Agilent E4980A LCR meter at
\SI{85}{\kilo\hertz}. Repeated readings varied within about
$\pm\SI{1.5}{\milli\ohm}$, and the absolute resistance drifted between sessions,
so each change is taken from the nominal resistance of the inductor alone
measured in the same session. \Cref{tab:app-rac-plate} gives the cold-plate
measurements of \cref{fig:rac_configurations}. Of the three changes, only the
\SI{1.0}{\milli\ohm} added by the Galinstan inside the plate lies within the
repeatability band.

\begin{table}[htbp]
  \centering
  \caption{AC resistance of the inductor on the cold plate.}
  \label{tab:app-rac-plate}
  \small
  \begin{tabular}{@{}lccc@{}}
    \toprule
    Configuration & $R_{\mathrm{ac}}$ & Change from & Change from \\
                  & (\si{\milli\ohm}) & inductor alone (\si{\milli\ohm})
      & empty plate (\si{\milli\ohm}) \\
    \midrule
    Inductor alone (nominal)  & 80.0  & --   & --   \\
    On empty cold plate       & 91.7  & 11.7 & --   \\
    On plate with Galinstan   & 92.7  & 12.7 & 1.0  \\
    On plate, vial beside     & 102.0 & 22.0 & 10.3 \\
    \bottomrule
  \end{tabular}
  \floatnote{Repeatability $\pm\SI{1.5}{\milli\ohm}$.}
\end{table}

\Cref{tab:app-rac-distance} gives the measurements with the vial of Galinstan
alone, without the cold plate, in a session whose nominal resistance was
\SI{82.15}{\milli\ohm}. The loss depends strongly on where the vial sits: it is
largest beside the ferrite core and not measurable beside the winding. Moving the vial away
from the side of the ferrite, the increase falls to within the repeatability band
by \SI{3}{\centi\metre}.

\begin{table}[htbp]
  \centering
  \caption{AC resistance with the Galinstan vial at different positions.}
  \label{tab:app-rac-distance}
  \small
  \begin{tabular}{@{}lcc@{}}
    \toprule
    Vial position & $R_{\mathrm{ac}}$ (\si{\milli\ohm}) & Change (\si{\milli\ohm}) \\
    \midrule
    None (nominal)              & 82.15 & --    \\
    Beside the winding          & 82.15 & 0     \\
    On top of the inductor      & 89.0  & 6.85  \\
    Beside the ferrite, right   & 99.0  & 16.85 \\
    Beside the ferrite, left    & 101.0 & 18.85 \\
    \addlinespace
    \SI{1}{\centi\metre} from the ferrite & 89.0  & 6.85 \\
    \SI{2}{\centi\metre} from the ferrite & 86.0  & 3.85 \\
    \SI{3}{\centi\metre} from the ferrite & 83.0  & 0.85 \\
    \SI{4}{\centi\metre} from the ferrite & 82.15 & 0    \\
    \SI{5}{\centi\metre} from the ferrite & 82.15 & 0    \\
    \bottomrule
  \end{tabular}
  \floatnote{No cold plate. Distances from the right side of the ferrite.
    Repeatability $\pm\SI{1.5}{\milli\ohm}$.}
\end{table}

\section{Flow and Coolant Temperatures with the Inductor Energised}

\Cref{tab:app-inductor-emf} lists the flow-sensor EMF as the inductor current was
raised in steps, with the pump held at \SI{18}{\ampere}. Without a filter, the
EMF fell by about \SI{3}{\percent}, from 151 to \SI{146.5}{\micro\volt}, as the
current rose to \SI{11.8}{\ampere} rms. With a filter on the sensor leads, the
readings were lower and rose irregularly, from 123 to as much as
\SI{151}{\micro\volt}, with no consistent trend. The change is therefore within the noise of the measurement, and whether
the field changes the flow is inconclusive (\cref{sec:results-inductor-thermal}).

\begin{table}[htbp]
  \centering
  \caption{Flow-sensor EMF against inductor current.}
  \label{tab:app-inductor-emf}
  \small
  \begin{tabular}{@{}cccc@{}}
    \toprule
    Inductor current & \multicolumn{2}{c}{Sensor EMF (\si{\micro\volt})}
      & Flow, no filter \\
    \cmidrule(lr){2-3}
    (\si{\ampere} rms) & No filter & With filter & (\si{\milli\litre\per\second}) \\
    \midrule
    0    & 151.0 & 123 & 1.39 \\
    2.0  & 150.0 & 123 & 1.38 \\
    4.0  & 151.0 & 125 & 1.39 \\
    5.9  & 152.0 & 126 & 1.40 \\
    7.9  & 150.0 & 137 & 1.38 \\
    9.9  & 148.5 & 151 & 1.37 \\
    11.8 & 146.5 & 136 & 1.35 \\
    \bottomrule
  \end{tabular}
  \floatnote{Pump at \SI{18}{\ampere}; flow from the calibration of
    \cref{tab:app-emf}.}
\end{table}

At steady state in the thermal test, the thermocouples read \SI{26.5}{\celsius}
at the cold-plate inlet and \SI{33.1}{\celsius} at the outlet, with the plate
surface at \SI{32.2}{\celsius} and ambient at \SI{21.8}{\celsius}. At
\SI{1.39}{\milli\litre\per\second}, the \SI{6.6}{\celsius} coolant temperature
rise corresponds to about \SI{21}{\watt} carried away by the Galinstan, roughly
half of the estimated \SI{42.76}{\watt} loss, with the rest lost from the upper
core and winding by natural convection. Because the temperature rise is only
about twice the $\pm\SI{3.1}{\celsius}$ uncertainty of a thermocouple difference
(\cref{app:uncertainty}), this estimate is indicative only, to about
$\pm\SI{10}{\watt}$.
